\documentclass[11pt]{article}
\usepackage[margin=1in]{geometry}
\usepackage{amsmath,amssymb,amsthm,booktabs}
\usepackage[hidelinks]{hyperref}
\usepackage{xurl}
\setlength{\emergencystretch}{2em}
\newtheorem{theorem}{Theorem}
\newtheorem{lemma}[theorem]{Lemma}
\newtheorem{corollary}[theorem]{Corollary}
\newcommand{\rk}{\operatorname{rank}}
\newcommand{\mr}{\operatorname{mr}}
\newcommand{\ar}{\operatorname{ar}}
\newcommand{\C}{\mathbb C}
\title{A generic rank-one obstruction for a seventeen-crossing knot}
\author{Written paired-span proof and exact signed controls}
\date{17 November 2026}
\begin{document}
\maketitle

Let $K_*$ be the closure of the four-braid
\[
 \beta=\sigma_1\sigma_2\sigma_3^{-2}\sigma_2
 \sigma_3^{-2}\sigma_2\sigma_1^2\sigma_2^3
 \sigma_3^{-1}\sigma_2^3 .
\]
Three local folds reduce its signed Wirtinger system to a previously
proved three-direction contradiction. The proof works over every
field and every nonzero meridian determinant, including unipotent
meridians. It uses roots and covectors separately and requires no
preserved form or invertible pairing matrix.
An explicit three-dimensional complex representation attains
augmentation rank three. With the independent meridional-rank
certificate this gives $\ar(K_*,\C)=3<4=\mr(K_*)=b(K_*)$.

For a finite-dimensional representation $\rho:G_{K_*}\to GL(V)$,
where $G_{K_*}=\pi_1(S^3\setminus K_*)$, define
\[
 S_\rho=\sum_{g\in G_{K_*}}\operatorname{im}(\rho(g)-I),\qquad
 r=\dim S_\rho,\qquad q=\operatorname{codim}_V V^{G_{K_*}}.
\]
The first sum is a span of vector displacement images.

\begin{theorem}\label{thm:pairedspans}
Over every field, every finite-dimensional representation of
$G_{K_*}$ with meridian displacement rank one satisfies
\[
 r\le4,\qquad q\le4,\qquad \min(r,q)\le3,\qquad r+q\le7.
\]
The representation need not be irreducible or completely reducible.
\end{theorem}

\section{The literal signed diagram}

Number PD ports cyclically $0,1,2,3$, with the understrand in
ports $0,2$. The following table gives the literal PD and its
signed Wirtinger rows. A row $(o,a,b,\varepsilon)$ means
\[
 M_b=M_o^\varepsilon M_aM_o^{-\varepsilon}.
\]
The understrand enters port $0$. The overstrand enters port $1$
when $\varepsilon=+$ and port $3$ when $\varepsilon=-$.
The segment labels at ports $0,2$ belong to arcs $a,b$, respectively;
both overports belong to arc $o$. Thus repeated PD segment labels
check the arc assignment directly.
\[
\begin{array}{c|rrrr|rrrr}
j&\multicolumn{4}{c|}{\mathrm{PD}_j}
 &\multicolumn{4}{c}{(o,a,b,\varepsilon)}\\\hline
0&4&2&5&1&0&1&2&-\\
1&5&17&6&16&8&2&3&-\\
2&27&6&28&7&3&13&14&+\\
3&7&28&8&29&14&3&4&+\\
4&17&9&18&8&4&8&9&-\\
5&29&18&30&19&9&14&15&+\\
6&19&30&20&31&15&9&10&+\\
7&9&21&10&20&10&4&5&-\\
8&2&22&3&21&10&0&1&-\\
9&22&4&23&3&1&10&11&-\\
10&23&11&24&10&5&11&12&-\\
11&11&25&12&24&12&5&6&-\\
12&25&13&26&12&6&12&13&-\\
13&31&26&32&27&13&15&16&+\\
14&13&33&14&32&16&6&7&-\\
15&33&15&0&14&7&16&0&-\\
16&15&1&16&0&0&7&8&-
\end{array}
\]
The exact signed binding to the closure of $\beta$ is retained in
\path{new17_generic_kch_inputs.json}; the optional topology replay
reconstructs this explicit braid, rather than consulting a knot name.

The seed arcs are $2,8,13,0$. Rows $1,\ldots,7,13$ give the new
arcs $3,14,4,9,15,10,5,16$ in that order. Reverse row $0$ to obtain
arc $1$, then rows $9,10,11,14$ give $11,12,6,7$.
These thirteen steps produce all seventeen meridians, proving that
the four seeds generate the knot group.

\section{Independent root and covector families}

Write a meridian image as
\[
 M_i=I-u_if_i,\qquad f_i(u_i)=d=1-\mu,\qquad
 M_i^{-1}=I+\mu^{-1}u_if_i.
\]
The roots $u_i$ and linear covectors $f_i$ are nonzero, and
$\mu=\det M_i\ne0$ is common to the conjugate meridians.
At a signed crossing propagate the pair by
\[
 (u_b,f_b)=(M_o^\varepsilon u_a,f_aM_o^{-\varepsilon}).
\]
Two pairs give the same nonzero rank-one displacement exactly when
they differ by $(u,f)\mapsto(cu,c^{-1}f)$, $c\ne0$.

Let $S,T$ be the spans of the four seed roots and covectors.
The seed matrices and their inverses preserve $S$.
Using $AB-I=(A-I)B+(B-I)$ shows $S=S_\rho$; also
$V^{G_{K_*}}=\bigcap_{\rm seeds}\ker f_i$.
Thus $r=\dim S\le4$ and $q=\dim T\le4$.
For a contradiction suppose $r=q=4$. Both seed families are
independent in their own spaces. We do not assume $S=V$, $T=V^*$
or a nonsingular evaluation matrix.

Put
\[
 A=M_2,\quad B=M_8,\quad C=M_{13},\quad D=M_0,
\]
and let $H$ be the span of the first three roots and $F$ the span
of their covectors. Write the pair of $D$ as $(v,\delta)$.
Then $v\notin H$ and $\delta\notin F$.
The eight early steps just listed use only $A,B,C$.
Consequently their pairs lie in $H,F$; in particular this holds for
\[
 X=M_{10},\qquad P=M_5,\qquad Y=M_{16}.
\]

\begin{lemma}[Three local folds]\label{lem:folds}
Under the simultaneous independence assumption $r=q=4$, the
literal relations force
\[
 X=A,\qquad M_{11}=D,\qquad P=C,\qquad M_6=D,\qquad Y=B.
\]
These deductions hold over every field and every $\mu\ne0$.
\end{lemma}
\begin{proof}
\emph{First fold.}
Rows $0,8$ give $DAD^{-1}=X^{-1}DX$.
Write the pairs of $A,X$ as $(a,f),(x,\xi)$ in $H,F$.
For some $\lambda\ne0$ equality of the rank-one pairs says
\[
\begin{split}
a-\delta(a)v&=\lambda\bigl(v+\mu^{-1}\xi(v)x\bigr),\\
f+\mu^{-1}f(v)\delta
 &=\lambda^{-1}\bigl(\delta-\delta(x)\xi\bigr).
\end{split}
\]
Comparing the fourth root and fourth covector components gives
$\lambda=-\delta(a)$ and $f(v)=\mu\lambda^{-1}$.
The remaining components are
\[
 a=\lambda\mu^{-1}\xi(v)x,\qquad
 f=-\lambda^{-1}\delta(x)\xi.
\]
Evaluation of the latter at $v$ gives $\delta(x)\xi(v)=-\mu$.
Multiplication of the two displayed identities therefore yields
$af=x\xi$, so $X=A$.
In particular $\delta(a)f(v)=-\mu$; no division by $d$ was used.
The group relation becomes $DAD=ADA$.
For $O=M_1=DAD^{-1}$ this gives $AO=OD$, so row $9$ gives
$M_{11}=O^{-1}AO=D$.

\emph{Second fold.}
Rows $10,11,12$ are the three inverse Hurwitz moves
\[
 (D,P)\longmapsto(P,N)\longmapsto(N,Q)\longmapsto(Q,C),
 \qquad N=M_{12},\quad Q=M_6.
\]
Choose $P=(p,\pi)$ in $H,F$ and set
$x=\pi(v)$, $y=\delta(p)$, $h=1+xy/\mu$.
Direct pair propagation gives
\[
\begin{array}{c|c|c}
 &\text{root}&\text{covector}\\\hline
N&v+(x/\mu)p&\delta-y\pi\\
Q&yv+hp&-(x/\mu)\delta+h\pi\\
C&hv+(x/\mu)(1+h)p&h\delta-y(1+h)\pi .
\end{array}
\]
The propagated pair on $C$ may differ from its seed pair by a
reciprocal nonzero rescaling, but its root still lies in $H$
and its covector in $F$. Either fourth component forces $h=0$.
Hence $xy=-\mu$, so $x,y$ are nonzero.
The last two pairs then have reciprocal scalar factors relative
to the pairs of $D,P$. Therefore $Q=D$ and $C=P$ as full matrices.

\emph{Third fold.}
Rows $14,15,16$ now give
\[
 (D,Y)\longmapsto(Y,M_7)\longmapsto(M_7,D)\longmapsto(D,B).
\]
Each inverse Hurwitz move preserves the ordered product.
Explicitly $DY=YM_7=M_7D=DB$, so $Y=B$.
Only powers of $\mu$ and scalars already proved nonzero have been
cancelled in these three folds.
\end{proof}

\section{The exact signed core}

\begin{lemma}\label{lem:coretransfer}
After the folds, the eight remaining signed relations are
equivalent to the $13a_{392}$ reduced system in
\cite[Lemma \texttt{lem:cores}]{extensions}. Their full root and
covector spans would both have dimension three.
\end{lemma}
\begin{proof}
The surviving relations are
\[
\begin{array}{rrrr}
(8,2,3,-),&(3,13,14,+),&(14,3,4,+),&(4,8,9,-),\\
(9,14,15,+),&(15,9,2,+),&(2,4,13,-),&(13,15,8,+).
\end{array}
\]
The bijection to the target arc labels, with $0=4$ and $6=7$
in that reduced system, is
\[
\begin{array}{c|rrrrrrrr}
\text{new}&2&3&4&8&9&13&14&15\\\hline
\text{target}&11&10&9&0&12&8&6&5 .
\end{array}
\]
The rows agree after solving some target rows for the other
under-arc, exchanging the two under-arcs and reversing the exponent
when needed. This is the same relation
$M_b=M_o^\varepsilon M_aM_o^{-\varepsilon}$ solved for $M_a$.
There is no parameter inversion.

Every new core pair lies in $H,F$, and the core contains the
independent original triples on arcs $2,8,13$.
Its full root and covector spans therefore both have dimension
three. The target's standard triple on arcs $0,6,12$ propagates
all eight core pairs. Its two spans must consequently also have
dimension three, making its roots and covectors independent
separately. This argument does not assume a nonsingular Gram
matrix or a nonzero minor of their evaluations.
\end{proof}

The precise imported conclusion of
\cite[Lemma \texttt{lem:cores}, $392$ case]{extensions} is:
for every field and every $\mu\ne0$, the full root and full
covector spans of this signed reduced eight-crossing system
cannot both have dimension three. Its proof uses the six
unrestricted directed pairings, with no old-clasp relation.
Lemma~\ref{lem:coretransfer} contradicts that conclusion.
Thus $r=q=4$ is impossible, proving
Theorem~\ref{thm:pairedspans}. In particular $r+q\le7$.
This proof also permits $\mu=1$, characteristic two and
primitive sixth-root parameters.

\section{Representation consequences}

\begin{corollary}\label{cor:irreducible}
Every irreducible finite-dimensional representation of $G_{K_*}$
over any field with meridian displacement rank one has dimension
at most three.
\end{corollary}
\begin{proof}
The nonzero invariant displacement span is all of $V$.
The common fixed subspace is zero: otherwise irreducibility
would make the representation trivial, contrary to the
nonzero meridian displacement. Hence $r=q=\dim V$ and
Theorem~\ref{thm:pairedspans} applies.
\end{proof}

\begin{corollary}\label{cor:forms}
If such a representation preserves a nondegenerate bilinear or
sesquilinear form over a field with involution, then $r=q\le3$,
even without complete reducibility. This includes symplectic
and Hermitian forms. A degenerate ambient form is not covered.
\end{corollary}
\begin{proof}
Take the form $B$ to be conjugate linear in its first argument
and linear in its second, with field involution $\sigma$;
the bilinear case has $\sigma=\mathrm{id}$.
Nondegeneracy defines adjoints by $B(Lx,y)=B(x,L^\dagger y)$.
Preservation gives $M_i^\dagger=M_i^{-1}$. For
$N_i=M_i-I=-u_if_i$,
\[
 N_i^\dagger=M_i^{-1}-I=-\mu^{-1}N_i,\qquad
 \operatorname{im}N_i^\dagger=ku_i .
\]
It follows that
$\ker f_i=\{x:B(x,u_i)=0\}$.
Thus $f_i$ is proportional to the linear covector
$x\mapsto\sigma(B(x,u_i))$.
The map $u\mapsto[x\mapsto\sigma(B(x,u))]$ is a semilinear
isomorphism, so the seed root and covector spans have equal
dimension. Apply Theorem~\ref{thm:pairedspans}.
\end{proof}

\begin{corollary}\label{cor:semisimple}
Every completely reducible finite-dimensional rank-one-meridian
representation has $r=q\le3$. In particular this holds for
finite-image representations when the image order is invertible
in the field, including every finite-image complex representation.
\end{corollary}
\begin{proof}
The ranks of the meridian displacements on irreducible summands
add to one. Exactly one summand is active. On each remaining
summand a meridian acts trivially, so normal generation makes
the entire knot group trivial there.
The active summand has dimension at most three by
Corollary~\ref{cor:irreducible}, giving $r=q\le3$.
If the finite image order is invertible, averaging a projection
onto any invariant subspace supplies an invariant complement;
induction gives complete reducibility.
\end{proof}

\begin{corollary}\label{cor:finiteclass}
For every finite complex image of $G_{K_*}$ with meridian
displacement rank one, its meridian conjugacy class generates
the image with at most three members.
\end{corollary}
\begin{proof}
The active summand in Corollary~\ref{cor:semisimple} is faithful
on the image and has dimension at most three.
The image on it is an irreducible finite complex reflection
group generated by the meridian class, by normal generation.
The imported Theorem \texttt{thm:lowrank} of \cite{finite}
states that any generating reflection class in an irreducible
finite complex reflection group of rank at most three generates
that group with at most three members.
\end{proof}

\section{Augmentation upper bound and explicit attainment}

We use Cornwell's augmentation rank $\ar(K,k)$ for field-valued
$U=1$ cord augmentations with $\epsilon(\mu)\ne1$,
with rank as in \cite[Definition~1.8]{cornwell}.

\begin{corollary}\label{cor:augmentation}
Over every field, each $U=1$ cord augmentation of $K_*$
with $\epsilon(\mu)\ne1$ has rank at most three.
The nonzero longitude value is unrestricted.
\end{corollary}
\begin{proof}
Cornwell's arbitrary-field augmentation-to-irreducible-KCH
correspondence and rank convention
\cite[Theorem~1.2 and Definition~1.8]{cornwell}
turn such an augmentation into an irreducible representation
whose dimension equals its rank.
Apply Corollary~\ref{cor:irreducible}.
This is a uniform argument over the admitted parameter values,
rather than a specialization from $\mu=-1$.
\end{proof}

\begin{theorem}\label{thm:attainment}
There is an irreducible three-dimensional complex KCH
representation of $G_{K_*}$ with meridian parameter $-1$.
\end{theorem}
\begin{proof}
Let $\phi=(1+\sqrt5)/2$, so $\phi^2=\phi+1$.
The following table specifies every arc root $u_i\in\mathbb Q(\phi)^3$.
Repeated arc indices use exactly the same vector.
\[
\begin{array}{c|rrr}
\text{arcs }i&\multicolumn{3}{c}{u_i}\\\hline
0,6,11&-1&\phi-1&-\phi\\
1&-2&0&0\\
2,10&-1&1-\phi&\phi\\
3,9&0&-2&0\\
4&-1&\phi-1&\phi\\
5,8,13,15,16&\phi-1&\phi&-1\\
7,12&-\phi&-1&1-\phi\\
14&\phi-1&-\phi&-1
\end{array}
\]
Every vector has squared norm four. Set
\[
 R_i=I-\tfrac12u_i u_i^T .
\]
Then $R_i^2=I$ and its eigenvalues are $-1,1,1$.
Substitution in the seventeen displayed signed rows gives
$R_b=R_o^\varepsilon R_aR_o^{-\varepsilon}$ in every case;
these are exact equalities in $\mathbb Q[\phi]/(\phi^2-\phi-1)$.
Equivalently $R_ou_a$ is $u_b$ up to sign.
The retained checker verifies all seventeen equations directly,
not just a reflection-class search.

The three roots on arcs $1,0,8$ have determinant $-4$.
Their pairings in that order are $2,2,2-2\phi$, all nonzero.
They form a basis with connected nonorthogonality graph.
For any nonzero invariant subspace, some one of these three
covectors is nonzero on it, since their common kernel is zero.
A displacement then puts its root into that subspace.
Conjugating by the other reflections propagates along the
nonorthogonality graph, so the subspace contains the full basis.
The representation is therefore irreducible.
No augmentation correspondence or classification of the finite
reflection image is needed for this representation assertion.
\end{proof}

\begin{corollary}\label{cor:finiteattainment}
The finite complex reflection-class lower bounds for $K_*$
have maximum exactly three. In particular the representation
of Theorem~\ref{thm:attainment} has finite image whose meridian
class requires exactly three generators.
\end{corollary}
\begin{proof}
The upper bound is Corollary~\ref{cor:finiteclass}.
For finiteness of the positive image, consider the thirty vectors
\[
 \mathcal R=\{\pm2e_i:1\le i\le3\}
 \ \cup\
 \{\text{cyclic permutations of }
     (\varepsilon_1,\varepsilon_2(\phi-1),\varepsilon_3\phi):
       \varepsilon_i=\pm1\}.
\]
The seed reflections permute this set, as verified exactly
using $\phi^2=\phi+1$. It spans $\C^3$, so the resulting
permutation action is faithful and the image is finite.
The four generating seed arcs $2,8,13,0$ give only three
distinct reflections because $u_8=u_{13}$.

Two members of the meridian class cannot generate this image.
Indeed, for two involutions $R,S$ in dimension three, choose an
eigenvector $w$ of $T=RS$ over $\C$.
Since $TR=RT^{-1}$, the nonzero space
$\langle w,Rw\rangle$ is invariant under $T,R$, and hence under
$S$, with dimension at most two. This contradicts the
irreducibility proved above. Thus the class rank of this finite
image is exactly three.
\end{proof}

\begin{corollary}\label{cor:gap}
\[
 \ar(K_*,\C)=3<4=\mr(K_*)=b(K_*).
\]
\end{corollary}
\begin{proof}
Cornwell's representation-to-augmentation construction
\cite[Proposition~2.12 and Definition~1.8]{cornwell}
applied to Theorem~\ref{thm:attainment} gives a rank-three
$U=1$ cord augmentation at $\mu=-1$.
Combine this with Corollary~\ref{cor:augmentation}.
The meridional-rank and bridge equalities are the separate
Theorem \texttt{thm:sharp} of \cite{new17};
they are not premises of the augmentation upper bound.
\end{proof}

\begin{corollary}\label{cor:sums}
For every integer $s\ge0$, put
$K_s=K_*\mathbin{\#}(\mathbin{\#}^{s}3_1)$. Then
\[
 \ar(K_s,\C)=s+3<s+4=\mr(K_s)=b(K_s).
\]
\end{corollary}
\begin{proof}
The representation in Theorem~\ref{thm:attainment} gives a seed
augmentation at $\mu=-1$ by the same source construction.
The trefoil has a rank-two augmentation at that same
parameter: the explicit irreducible two-dimensional matrices
$\left(\begin{smallmatrix}-1&1\\0&1\end{smallmatrix}\right)$ and
$\left(\begin{smallmatrix}1&0\\1&-1\end{smallmatrix}\right)$
satisfy the trefoil relation and are retained in
\cite[the connected-sum corollary and its control]{generic373}.
Denote these two matrices by $A,B$. The matrices $I,A,B,AB$
form a basis of $M_2(\C)$, giving irreducibility directly.
The two meridians generate the trefoil group, so the same
seed-span argument bounds every irreducible KCH dimension
there by two. Its augmentation rank is therefore exactly two.
Cornwell's Theorem~1.9 gives a common-parameter augmentation of
rank $d_1+d_2-1$ on a connected sum.
The upper part of that theorem starts with an arbitrary
sum augmentation and bounds its rank by
$\ar(K_1,\C)+\ar(K_2,\C)-1$, without fixing its parameter.
Iterating both parts yields $\ar(K_s,\C)=s+3$.
The independent meridional and bridge formula is
\cite[Corollary \texttt{cor:trefoils}]{new17}.
No general meridional-rank additivity is assumed.
\end{proof}

Strict augmentation-rank gaps were already known: Cornwell's
Corollary~1.12 gives $2=\ar(K,\C)<\mr(K)=3$ for
$8_{16},8_{17},10_{91},10_{94}$.
No first-counterexample or worldwide-priority assertion is made.

\section{Exact controls and scope}

The default \path{verify_new17_generic_kch.py} replay uses exact
arithmetic. It binds the literal PD, signed rows and four-seed trace,
checks the folds and signed core bijection, and verifies the
three-dimensional positive representation and its faithful
action on the thirty-vector set.
Its input is \path{new17_generic_kch_inputs.json}; the output
receipt is \path{new17_generic_kch_verification.json}.
The optional \texttt{--topology} mode reconstructs the explicit
four-braid with the pinned topology software.
The retained positive and corrupted controls test actual diagram
equations rather than only a formal matrix quartet.

The reduced $392$ contradiction is an explicit written and exact
algebra dependency, separate from the new diagram binding.
The representation theorem permits arbitrary finite ambient
dimension, singular evaluation matrices, every field and every
$\mu\ne0$. Its unrestricted conclusion is
$\min(r,q)\le3$ and $r+q\le7$.
Separate bounds $r,q\le3$ use the hypotheses of the form,
irreducibility or complete-reducibility corollaries.
The augmentation correspondence excludes $\epsilon(\mu)=1$.
These are internal written proofs and exact replays, not a
claim of formal verification or external specialist review.

\begin{thebibliography}{9}\small
\bibitem{extensions} \emph{Signed generic rank-one obstructions for
sixteen knots}, companion \path{GENERIC_KCH_EXTENSIONS.tex},
Lemma \texttt{lem:cores}, $392$ case.
\bibitem{cornwell} C.~R.~Cornwell,
\emph{KCH representations, augmentations, and $A$-polynomials},
J.\ Symplectic Geom.\ 15 (2017), 983--1017.
Retained arXiv:1310.7526v3, Theorem~1.2, Definition~1.8,
Proposition~2.12, Theorem~1.9 and Corollary~1.12.
\bibitem{new17} \emph{An explicit seventeen-crossing four-bridge knot},
companion \path{NEW_SEVENTEEN_CROSSING_KNOT.tex},
Definition \texttt{def:knot}, Theorem \texttt{thm:sharp},
Corollary \texttt{cor:trefoils}.
\bibitem{finite} \emph{A finite complex reflection boundary for
sixteen knots}, companion
\path{FINITE_COMPLEX_REFLECTION_BOUNDARY.tex},
Theorem \texttt{thm:lowrank}. Its classification premise is the
retained primary Shephard--Todd theorem.
\bibitem{generic373} \emph{A generic rank-one obstruction for
$13a_{373}$}, companion \path{GENERIC_KCH_OBSTRUCTION.tex},
Corollary \texttt{cor:sums} and the exact trefoil control.
\end{thebibliography}
\end{document}
